\chapter{Non-Functional Requirements}
\label{chapter:requirements}

The following non-functional requirements (NFRs) define constraints on how the integrated solution should perform and behave. They're based on enterprise infrastructure best practices, industry standards (CIS Kubernetes Benchmark, NIST SP~800-207), and the project context.

\section{Scalability}

\begin{itemize}
    \item \textbf{NFR-SC-01:} The monitoring stack must work smoothly on clusters with up to 100 worker nodes and 500 running pods without needing manual reconfiguration.
    \item \textbf{NFR-SC-02:} The AI inference pipeline must scale horizontally — throughput should increase linearly when we add more inference replicas.
    \item \textbf{NFR-SC-03:} Cilium must operate at wire speed on 10~Gbps network interfaces without dropping packets under normal cluster load.
\end{itemize}

\section{Performance}

\begin{itemize}
    \item \textbf{NFR-PE-01:} The extra CPU overhead from Cilium (eBPF programs) must not exceed 5\% of total cluster CPU capacity compared to a Flannel baseline under the same workloads.
    \item \textbf{NFR-PE-02:} The extra memory footprint of the Cilium agent must not exceed 200~MiB per node.
    \item \textbf{NFR-PE-03:} End-to-end monitoring latency (from when a metric is generated to when an alert is sent) must be under 60 seconds for critical severity alerts.
    \item \textbf{NFR-PE-04:} AI model inference latency for a single prediction must be below 500~ms at the 99th percentile.
\end{itemize}

\section{Security}

\begin{itemize}
    \item \textbf{NFR-SE-01:} All communication between services in the cluster must be controlled by explicit CiliumNetworkPolicy resources. We'll enforce a default-deny approach for both incoming and outgoing traffic.
    \item \textbf{NFR-SE-02:} Mutual TLS (mTLS) must be automatically enforced between all workloads using Cilium's WireGuard or IPsec encryption, without requiring any application code changes.
    \item \textbf{NFR-SE-03:} The solution must pass the applicable controls of the CIS Kubernetes Benchmark~v1.9.
    \item \textbf{NFR-SE-04:} All container images used in the monitoring and security stack must be scanned for known vulnerabilities before deployment. We won't deploy images with critical CVEs.
    \item \textbf{NFR-SE-05:} Secrets (API keys, certificates, credentials) must not be stored in plaintext. They'll be managed via Kubernetes Secrets with RBAC access controls or an external secrets manager.
\end{itemize}

\section{Reliability}

\begin{itemize}
    \item \textbf{NFR-RE-01:} The monitoring stack (Prometheus, Grafana, AlertManager) must be deployed in high-availability mode with at least two replicas per component and anti-affinity rules to prevent them from running on the same node.
    \item \textbf{NFR-RE-02:} We'll retain metric data for at least 30 days, with long-term storage backed by persistent volumes.
    \item \textbf{NFR-RE-03:} The Cilium agent must recover automatically from pod restarts without manual intervention, re-establishing all network policies within 30 seconds.
    \item \textbf{NFR-RE-04:} The AI anomaly detection service must be resilient to losing a single input data stream. It should keep producing predictions using available data and send a degraded-mode alert.
\end{itemize}

\section{Observability}

\begin{itemize}
    \item \textbf{NFR-OB-01:} Full Layer~3 through Layer~7 network flow visibility must be available via Hubble for all traffic within the cluster and going out of the cluster.
    \item \textbf{NFR-OB-02:} The solution must support the three pillars of observability: metrics (Prometheus), logs (Loki), and distributed traces (OpenTelemetry/Jaeger), with the ability to correlate across all three.
    \item \textbf{NFR-OB-03:} Grafana dashboards must provide at minimum: cluster-level resource usage, per-namespace network flow heatmaps, anomaly detection score time series, and counts of active network policy violations.
\end{itemize}

\section{Maintainability and Portability}

\begin{itemize}
    \item \textbf{NFR-MA-01:} All components must be deployed exclusively via Helm charts stored in a version-controlled Git repository. No manual \texttt{kubectl apply} steps should be needed after initial cluster setup.
    \item \textbf{NFR-MA-02:} The GitOps reconciliation loop (Flux CD) must automatically bring the cluster to the desired state defined in Git within 5 minutes of a commit.
    \item \textbf{NFR-MA-03:} The solution must be validated on at least two different Kubernetes distributions: bare-metal (kubeadm) and a managed cloud service (AKS or EKS).
\end{itemize}